%%---Main Packages-----------------------------------------------------------------------
\usepackage[english, ngerman]{babel}	%Mul­tilin­gual sup­port for LaTeX
\addto\extrasngerman{%
  \renewcommand{\sectionautorefname}{Abschnitt}%
  \renewcommand{\subsectionautorefname}{Unterabschnitt}%
  \renewcommand{\subsubsectionautorefname}{Unterunterabschnitt}%
  \renewcommand{\tableautorefname}{Tabelle}%
}
\usepackage[T1]{fontenc}				      %Stan­dard pack­age for se­lect­ing font en­cod­ings
\usepackage[utf8]{inputenc}				  %Ac­cept dif­fer­ent in­put en­cod­ings
\usepackage{lmodern}                 %The newer Font-Set
\usepackage{textcomp}					      %LaTeX sup­port for the Text Com­pan­ion fonts
\usepackage{graphicx} 					      %En­hanced sup­port for graph­ics
\usepackage{svg}
\usepackage{float}						        %Im­proved in­ter­face for float­ing ob­jects
%\usepackage{ifdraft}                %Let you check if the doc is in draft mode

%%---Useful Packages---------------------------------------------------------------------
% \usepackage[pdftex,dvipsnames]{xcolor}  %Driver-in­de­pen­dent color ex­ten­sions for LaTeX
\usepackage{csquotes}                   %Simpler quoting with \enquote{}
\usepackage{siunitx} 					     %A com­pre­hen­sive (SI) units pack­age
\usepackage{listings}					     %Type­set source code list­ings us­ing LaTeX
\lstset{
  basicstyle=\ttfamily\small,
  breaklines=true,
  breakatwhitespace=false,  % allow breaking inside long tokens (URLs, keys, etc.)
  columns=fullflexible,      % let listings adjust character widths to enable breaks
  keepspaces=true,
  showstringspaces=false,
  tabsize=2,
  % optional: visual cue for wrapped lines
  postbreak=\mbox{\textcolor{gray}{$\hookrightarrow$}\space}
}
\usepackage[bottom]{footmisc}			  %A range of foot­note op­tions
\usepackage{footnote}					     %Im­prove on LaTeX's foot­note han­dling
\usepackage{verbatim}					     %Reim­ple­men­ta­tion of and ex­ten­sions to LaTeX ver­ba­tim
\usepackage[textsize=footnotesize]{todonotes} %Mark­ing things to do in a LaTeX doc­u­ment
\usepackage{lipsum}              % Gives you access to blindtext
\usepackage[acronym]{glossaries}
\makeglossaries 
\usepackage{pgfgantt}
\hypersetup{colorlinks=true, urlcolor=blue, linkcolor=black, citecolor=blue}
\usepackage{soul}
\usepackage{msc} % For sequence diagrams
\usepackage{makecell}
\usepackage{changepage}
%%---Tikz Packages-----------------------------------------------------------------------
%\usepackage{standalone}
\usepackage{tikz}
\usetikzlibrary{shapes.geometric, arrows.meta, positioning}
\tikzstyle{block} = [
    rectangle,
    rounded corners,
    minimum height=1.2cm,
    % minimum width=3cm,
    text centered,
    draw=black,
    fill=blue!20,
    align=center
]
\tikzstyle{arrow} = [thick,->,>=stealth]
% \usepackage{tikz-uml}
\usetikzlibrary{shapes.geometric, arrows.meta, positioning}
%\usepackage{circuitikz}
%\usetikzlibrary{arrows}
%\usetikzlibrary{calc}
%\usetikzlibrary{intersections}

%%---Math Packages-----------------------------------------------------------------------
\usepackage{amsmath}					    %AMS math­e­mat­i­cal fa­cil­i­ties for LaTeX
%\usepackage{amssymb}					  %Type­set­ting symbols (AMS style)
\usepackage{ragged2e,array}						  %Ex­tend­ing the ar­ray and tab­u­lar en­vi­ron­ments
\newcolumntype{Y}{>{\RaggedRight\arraybackslash}X} % ragged-right wrapping
%\usepackage{amsthm}					    %Type­set­ting the­o­rems (AMS style)

%%---Table Packages----------------------------------------------------------------------
\usepackage{tabularx}					  %Tab­u­lars with ad­justable-width columns
%\usepackage{longtable}
\usepackage{ltablex}   % combines longtable + tabularx
\keepXColumns   
\usepackage{multirow}					  %Create tab­u­lar cells span­ning mul­ti­ple rows
\usepackage{multicol}					  %In­ter­mix sin­gle and mul­ti­ple columns

%%---PDF / Figure Packages---------------------------------------------------------------
\usepackage{pdfpages}					  %In­clude PDF doc­u­ments in LaTeX
\usepackage{pdflscape}					  %Make land­scape pages dis­play as land­scape
%\usepackage{subfig}					    %Fig­ures di­vided into sub­fig­ures

%%---Other Packages----------------------------------------------------------------------
%\usepackage{xargs}              %De­fine com­mands with many op­tional ar­gu­ments
%\usepackage{drawio}              %In­clude draw.io di­a­grams di­rectly in LaTeX

%%---Main Settings-----------------------------------------------------------------------
\graphicspath{{./graphics/}}			%Defines the graphicspath
%\geometry{twoside=false}				%twoside=false disables the "bookstyle"
\setlength{\marginparwidth}{2cm}
\overfullrule=5em						    %Creates a black rule if text goes over the margins => debugging

%%---User Definitions--------------------------------------------------------------------
%%Tabel-Definitions: (requires \usepackage{tabularx})
\newcolumntype{L}[1]{>{\raggedright\arraybackslash\sloppy\hspace{0pt}}p{#1}}    %column-width and alignment
\newcolumntype{C}[1]{>{\centering\arraybackslash}p{#1}}
\newcolumntype{R}[1]{>{\raggedleft\arraybackslash}p{#1}}					

% \newcommand{\ac}[1]{\gls{#1}}
\newcommand{\ac}[1]{%
  \ifglsused{#1}%
    {\gls{#1}\textsuperscript{*}}% If already used → add asterisk
    {\gls{#1}}% If first time → no asterisk
}
% Set caption labels for figures and their references
\renewcaptionname{ngerman}{\figurename}{Abbildung}
\addto\extrasngerman{\renewcommand{\figureautorefname}{Abbildung}}

% Command for referencing sections with number and name
% \newcommand*{\fullref}[1]{\hyperref[{#1}]{\ref*{#1} \nameref*{#1}}}
\newcommand*{\fullref}[1]{%
  \autoref{#1} (\nameref{#1})%
}